\documentclass[12pt]{article}
\usepackage{amsmath, amssymb, amsthm}
\usepackage{geometry}
\geometry{a4paper, margin=2.5cm}

\title{p-adische Zahlen als Vektoren: Skala, Richtung und Geometrie}
\author{}
\date{}

\begin{document}
\maketitle

\section{Grundidee}

Eine p-adische Zahl kann sinnvoll als Vektor interpretiert werden.
Dabei ist \emph{Vektor} nicht im euklidischen Sinn zu verstehen, sondern
als unendliche Koordinatenstruktur relativ zur Basis \(p\).

Ziel dieses Textes ist es, die folgenden beiden Aussagen formal einzuordnen:
\begin{itemize}
  \item Die \textbf{Länge} eines p-adischen Vektors ergibt sich aus der Division durch Potenzen von \(p\).
  \item Die \textbf{Richtung} (umgangssprachlich: Winkel) ergibt sich aus der Modulo-Struktur.
\end{itemize}

\section{p-adische Darstellung als Koordinatenvektor}

Sei \(p\) eine Primzahl. Jede p-adische Zahl \(x \in \mathbb{Q}_p\) besitzt eine eindeutige Darstellung
\[
x = \sum_{k = k_0}^{\infty} a_k p^k,
\quad a_k \in \{0,1,\dots,p-1\}, \quad k_0 \in \mathbb{Z}.
\]

Diese Darstellung kann als unendlicher Koordinatenvektor gelesen werden:
\[
x \;\longleftrightarrow\; (a_{k_0}, a_{k_0+1}, a_{k_0+2}, \dots).
\]

Die Basis ist dabei nicht orthonormal, sondern hierarchisch skaliert durch die Gewichte \(p^k\).

\section{Länge: Bewertung und Division}

Die p-adische Bewertung \(v_p : \mathbb{Q}_p^\times \to \mathbb{Z}\) ist definiert durch
\[
x = p^{v_p(x)} \cdot u,
\quad u \in \mathbb{Z}_p^\times.
\]

Der zugehörige p-adische Betrag ist
\[
|x|_p = p^{-v_p(x)}.
\]

\subsection*{Interpretation}

\begin{itemize}
  \item Die \emph{Division} durch Potenzen von \(p\) isoliert die Skala eines Elements.
  \item Die Länge des p-adischen Vektors ist vollständig durch \(v_p(x)\) bestimmt.
\end{itemize}

Damit gilt:

\[
\textbf{Länge} \;\equiv\; \text{Bewertung} \;\equiv\; \text{Division durch } p^{v_p(x)}.
\]

Dies ist mathematisch exakt und nicht metaphorisch.

\section{Richtung: Modulo-Struktur}

Nach Herausziehen der Skala bleibt der Einheitsfaktor
\[
u = x \cdot p^{-v_p(x)} \in \mathbb{Z}_p^\times.
\]

Die Einheitengruppe zerfällt strukturell als
\[
\mathbb{Z}_p^\times \cong \mu_{p-1} \times (1 + p\mathbb{Z}_p),
\]
wobei
\begin{itemize}
  \item \(\mu_{p-1}\) die endliche zyklische Gruppe der Einheiten modulo \(p\) ist,
  \item \(1 + p\mathbb{Z}_p\) den infinitesimalen Richtungsanteil enthält.
\end{itemize}

\subsection*{Modulo als Richtungsinformation}

Die Reduktion
\[
x \bmod p^n
\]
fixiert die ersten \(n\) p-adischen Ziffern und damit eine Richtung bis zu einer gegebenen Tiefe.

\begin{itemize}
  \item Modulo \(p\): grobe Richtung (diskrete Phase)
  \item Modulo \(p^n\): zunehmend feinere Richtungsauflösung
\end{itemize}

Es existiert kein Winkel im euklidischen Sinn, da:
\begin{itemize}
  \item kein Skalarprodukt existiert,
  \item die Ultrametrik keine Winkel zulässt.
\end{itemize}

Dennoch übernimmt die Modulo-Struktur exakt die Rolle einer \emph{Orientierung}.

\section{Ultrametrische Geometrie}

Der p-adische Abstand ist definiert durch
\[
d(x,y) = |x-y|_p.
\]

Er erfüllt die starke Dreiecksungleichung
\[
|x+y|_p \le \max(|x|_p, |y|_p).
\]

Konsequenzen:
\begin{itemize}
  \item Kugeln sind verschachtelt oder disjunkt.
  \item Nähe bedeutet: gleiche Anfangskoordinaten.
  \item Die Geometrie ist baumartig, nicht linear.
\end{itemize}

\section{Zusammenfassung}

Eine p-adische Zahl zerfällt eindeutig in
\[
x = \underbrace{p^{v_p(x)}}_{\text{Skala / Länge}}
\cdot
\underbrace{u}_{\text{Richtung / Modulo-Struktur}}.
\]

Damit gilt:

\[
\boxed{
\begin{aligned}
\text{Länge} &\;\longleftrightarrow\; \text{Division (Bewertung)} \\
\text{Richtung} &\;\longleftrightarrow\; \text{Modulo-Struktur}
\end{aligned}
}
\]

Oder in einem Satz:

\begin{center}
\emph{In der p-adischen Geometrie liefert die Division die Skala, und das Modulo die Orientierung.}
\end{center}

Diese Interpretation ist strukturell korrekt und ersetzt Winkel und Normen durch Bewertung und Restklassen.

\end{document}

